\section{Introduction}

\IEEEPARstart{R}{obotic} manipulation is fundamentally grounded in three-dimensional geometry. A policy must determine where an object can be grasped, which region should make contact, and how the end effector or a tool should move relative to the object. These decisions often depend on functional parts rather than object identity alone: a mug handle supports grasping and hanging, a hammer handle and head serve different roles, and the opening and spout of a container constrain stirring and pouring. Because the geometry of these parts varies across instances, successful manipulation requires a representation that preserves spatial structure while exposing functional correspondence across objects.

Point clouds provide a natural 3D interface for this purpose and have enabled effective visuomotor policies~\cite{ze20243d,ke20243d}. In practice, however, a point cloud is not the object surface, but a partial measurement produced by a particular camera arrangement, viewpoint, depth sensor, calibration and sampling procedure. Even when the physical scene is unchanged, the observed point set can vary because of self-occlusion, missing depth, reflective materials, mask errors, and random resampling. These effects become more pronounced in real-world deployment than in simulation. A policy trained directly on such observations must therefore learn the manipulation mapping while simultaneously resolving functional-part ambiguity and suppressing observation-specific variation, often from a limited number of demonstrations.

\begin{figure*}[t]
    \vspace{-15pt}
    \centering
    \includegraphics[width=1\linewidth]{Images/teaser.png}
    \vspace{1pt}
    \caption{Teaser of our object-centric continuous semantic field. The point-attached 2D and 3D baselines attach semantics to observation-dependent samples. We instead condition a queryable semantic field on the object point cloud and read part-aware embeddings at explicit 3D query locations, producing semantic point clouds for downstream policy learning.}
    \label{fig:teaser}
\end{figure*}

A natural response is to enrich observed geometry with semantic features, by lifting 2D foundation-model descriptors into 3D~\cite{chen2025g3flow,fan2026any3d,wang2024gendp} or using point-wise 3D semantic encoders~\cite{chen2026learning,xu2026action}. Such features expose functional regions, but semantic enrichment does not automatically remove observation dependence. Coordinates and their attached descriptors can still change together, requiring the policy to resolve correspondences across both object instances and measurements of the same instance.


This limitation motivates separating observation from semantic readout. We use the observed object cloud as a \emph{support set} conditioning a continuous field and evaluate embeddings at independently resampled \emph{query locations}. Queries remain observation-dependent rather than fixed canonical points. The decoupling lies in decoding each embedding from the complete support-conditioned field instead of attaching it directly to a support sample. Training under perturbed observations encourages corresponding object regions to retain similar semantic responses.


We instantiate this idea as a category-level semantic field learned with part supervision. A frozen 3D encoder provides support features for an object-conditioned tri-plane cache, and a coordinate-conditioned decoder produces an embedding at each query. Part supervision identifies functional regions, cross-instance alignment associates corresponding parts, and consistency regularization encourages reduced sensitivity to changes in the support observation. After training, the frozen field exports query coordinates and embeddings as a semantic point cloud for a DP3-style policy. This interface preserves explicit 3D structure and adds part-aware information without changing the policy objective or action representation.


We evaluate the method from both representation and manipulation perspectives. Representation analysis examines whether the learned field produces consistent part-aware features across object observations and instances, while ablation experiments assess the contribution of its main components and training objectives. We measure downstream policy performance on four RoboTwin simulation tasks~\cite{chen2025robotwin} and four real-world bimanual manipulation tasks. Together, these experiments test the quality of the learned representation and its practical value for policy learning. Our contributions are threefold:
\begin{itemize}
    \item We formulate a support/query-decoupled object representation that reads functional-part features through an object-conditioned continuous field instead of attaching them directly to noisy support samples. The formulation preserves a point-cloud interface while allowing semantic readout at resampled 3D locations.
    \item We learn category-level fields using part anchoring, cross-instance alignment, and support-perturbation consistency, and qualitatively examine part discrimination and cross-instance consistency. Simulation ablations assess the three learning objectives and the query sampling design.
    \item We integrate the resulting semantic point clouds into a DP3-based policy with an unchanged action objective and evaluate manipulation in simulation and on a real bimanual platform, including held-out real-world object instances.
\end{itemize}
